\chapter[Migraciones en Producción]{\emj{🔧}~Migraciones en Producción}

\section[El problema de las migraciones en producción]{\emj{⚠️}~El problema de las migraciones en producción}

Alterar una tabla con millones de filas en producción sin downtime es uno de los desafíos más subestimados del DBA Senior.
Un \texttt{ALTER TABLE ADD COLUMN NOT NULL} en una tabla de 100M filas puede bloquear todas las escrituras por minutos.

\begin{warningbox}
\textbf{Locks en DDL:} \texttt{ALTER TABLE} toma un \texttt{AccessExclusiveLock} que bloquea \emph{todas} las lecturas y escrituras mientras dura.
Si hay una transacción larga corriendo, el ALTER espera. Y mientras espera, \emph{las queries nuevas también esperan} (cola de locks).
Un ALTER de 5 segundos puede provocar 500 queries encoladas que provocan timeout cascada y aparente caída total del servicio.
\end{warningbox}

\section[Expand-Contract Pattern]{\emj{📋}~Expand-Contract Pattern}

El patrón para migraciones sin downtime. Divide cada cambio en tres fases deployables independientemente.

\begin{teoriabox}{Las 3 fases del Expand-Contract}
\begin{enumerate}
  \item \textbf{Expand (Add):} agregar la nueva estructura (columna, tabla, índice) sin eliminar la vieja. El código nuevo puede leer ambas; el código viejo sigue funcionando solo con la vieja.
  \item \textbf{Migrate (Backfill):} llenar la nueva estructura con datos existentes. Opcionalmente: actualizar el código para escribir en ambas.
  \item \textbf{Contract (Remove):} eliminar la estructura vieja una vez que todo el código usa la nueva.
\end{enumerate}
\end{teoriabox}

\begin{lstlisting}[caption={Expand-Contract: renombrar columna (ejemplo real)}]
-- OBJETIVO: renombrar orders.user_id → orders.customer_id
-- Esto NO se hace con: ALTER TABLE orders RENAME COLUMN user_id TO customer_id
-- (rompería todo el código viejo instantáneamente)

-- FASE 1 - EXPAND: agregar nueva columna
ALTER TABLE orders ADD COLUMN customer_id BIGINT;
-- Agregar índice (concurrente, no bloquea):
CREATE INDEX CONCURRENTLY idx_orders_customer_id ON orders (customer_id);

-- FASE 2 - BACKFILL: copiar datos en batches (NO en una sola query)
-- Script de backfill por lotes:
DO $$
DECLARE
  batch_size  INT := 5000;
  last_id     BIGINT := 0;
  max_id      BIGINT;
BEGIN
  SELECT MAX(id) INTO max_id FROM orders;
  WHILE last_id < max_id LOOP
    UPDATE orders
    SET customer_id = user_id
    WHERE id > last_id AND id <= last_id + batch_size
      AND customer_id IS NULL;
    last_id := last_id + batch_size;
    PERFORM pg_sleep(0.1);  -- throttle para no saturar I/O
  END LOOP;
END $$;

-- Deploy del código que escribe en AMBAS columnas (user_id Y customer_id)

-- FASE 3 - CONTRACT: remover columna vieja
-- Primero asegurar que customer_id tiene NOT NULL (una vez backfill completo):
ALTER TABLE orders ALTER COLUMN customer_id SET NOT NULL;
-- Luego eliminar la vieja:
ALTER TABLE orders DROP COLUMN user_id;
-- Eliminar índice viejo si existía:
DROP INDEX CONCURRENTLY idx_orders_user_id;
\end{lstlisting}

\section[CREATE INDEX CONCURRENTLY]{\emj{🔄}~CREATE INDEX CONCURRENTLY}

\begin{lstlisting}[caption={Diferencias entre index normal y CONCURRENTLY}]
-- Index normal: bloquea escrituras mientras construye
-- Tarda segundos a minutos dependiendo del tamaño
CREATE INDEX idx_orders_status ON orders (status);

-- CONCURRENTLY: no bloquea escrituras (solo reads afectados brevemente al final)
-- Tarda más (2-3 pasadas), pero es seguro en producción
CREATE INDEX CONCURRENTLY idx_orders_status ON orders (status);

-- Si falla (crash, cancelación): queda en estado INVALID
-- Verificar:
SELECT indexname, indisvalid
FROM pg_indexes JOIN pg_index ON indexrelid = (
  SELECT oid FROM pg_class WHERE relname = indexname
)
WHERE tablename = 'orders';

-- Limpiar index inválido:
DROP INDEX CONCURRENTLY idx_orders_status;
-- Luego recrear
\end{lstlisting}

\section[pg\_repack: reorganizar sin bloquear]{\emj{🗃️}~pg\_repack: reorganizar sin bloquear}

\texttt{VACUUM FULL} compacta una tabla pero la bloquea completamente.
\texttt{pg\_repack} hace lo mismo construyendo la tabla nueva en segundo plano y haciendo swap atómico al final.

\begin{lstlisting}[language=, caption={Uso de pg\_repack}]
# Instalar extensión (en la BD):
# CREATE EXTENSION pg_repack;

# Reorganizar tabla (equivalente a VACUUM FULL sin lock prolongado):
pg_repack -h localhost -U postgres -d mydb -t orders

# Reorganizar tabla con nuevo orden físico (CLUSTER equivalente):
pg_repack -h localhost -U postgres -d mydb -t orders --order-by created_at

# Solo índices (reorganizar bloat en índices):
pg_repack -h localhost -U postgres -d mydb -t orders --only-indexes

# Toda la base de datos:
pg_repack -h localhost -U postgres -d mydb
\end{lstlisting}

\section[Flyway y Liquibase]{\emj{🗄️}~Flyway y Liquibase}

Herramientas de versionado de esquema: registran qué migraciones se han aplicado y aplican las pendientes en orden.

\begin{lstlisting}[language=, caption={Convenciones de Flyway}]
# Estructura de archivos:
# db/migration/
#   V1__init_schema.sql
#   V2__add_orders_table.sql
#   V3__add_customer_index.sql
#   R__create_reports_view.sql   (repeatable: se reaplica si cambia)

# Naming: V{version}__{description}.sql
# Version puede ser numérica (1, 2, 3) o timestamp (20250115120000)

# Aplicar migraciones pendientes:
flyway migrate

# Ver estado:
flyway info

# Validar que migraciones aplicadas coinciden con archivos:
flyway validate

# Reparar tabla de migraciones (si hay migration fallida):
flyway repair
\end{lstlisting}

\begin{lstlisting}[language=, caption={Liquibase con changesets XML}]
<!-- db/changelog/changelog.xml -->
<!-- <databaseChangeLog>
  <changeSet id="1" author="aldo">
    <createTable tableName="customers">
      <column name="id" type="BIGINT" autoIncrement="true">
        <constraints primaryKey="true"/>
      </column>
      <column name="email" type="VARCHAR(255)">
        <constraints nullable="false" unique="true"/>
      </column>
    </createTable>
  </changeSet>

  <changeSet id="2" author="aldo">
    <addColumn tableName="orders">
      <column name="customer_id" type="BIGINT"/>
    </addColumn>
  </changeSet>
</databaseChangeLog> -->

# Aplicar:
liquibase update

# Rollback al changeset anterior:
liquibase rollback --tag=v1.0
\end{lstlisting}

\begin{lstlisting}[caption={Agregar NOT NULL sin downtime}]
-- Problema: ALTER TABLE t ALTER COLUMN c SET NOT NULL
-- escanea toda la tabla para verificar que no hay NULLs → lock largo

-- Solución moderna (PG 12+):
-- 1. Agregar constraint de check primero (sin NOT VALID → no escanea):
ALTER TABLE orders ADD CONSTRAINT orders_customer_id_not_null
  CHECK (customer_id IS NOT NULL) NOT VALID;

-- 2. Validar en background (solo ShareUpdateExclusiveLock, no bloquea):
ALTER TABLE orders VALIDATE CONSTRAINT orders_customer_id_not_null;

-- 3. Ahora PG sabe que no hay NULLs → el SET NOT NULL es instantáneo:
ALTER TABLE orders ALTER COLUMN customer_id SET NOT NULL;

-- 4. Eliminar el constraint temporal:
ALTER TABLE orders DROP CONSTRAINT orders_customer_id_not_null;
\end{lstlisting}

\begin{entrevistabox}
\begin{enumerate}
  \item Explica el Expand-Contract pattern. ¿Por qué no puedes simplemente renombrar una columna en un ALTER directo?
  \item ¿Cuál es la diferencia entre \texttt{CREATE INDEX} y \texttt{CREATE INDEX CONCURRENTLY}? ¿Qué pasa si falla?
  \item ¿Cómo agregas una columna NOT NULL a una tabla de 500M filas sin downtime?
  \item ¿Qué pasa si una migración de Flyway falla a mitad? ¿Cómo lo resuelves?
  \item ¿Cuándo usarías pg\_repack vs VACUUM FULL?
  \item ¿Qué es un ``lock queue'' en PostgreSQL y cómo un DDL puede tumbar tu servicio aunque tarde solo 2 segundos?
\end{enumerate}
\end{entrevistabox}
